% !PRINTED 416
\weberpage{416}
Ist $p' > 1$, also $\pi > 1$, so fehlt in dieser Gleichung das
Glied mit der $(\nu-1)^{\text{ten}}$ Potenz der Unbekannten, dessen Co\"efficient der negativen Summe der Wurzeln gleich ist. Wir haben
also den Satz:

\gesperrt{Die Summe aller primitiven $n^{\text{ten}}$ Einheitswurzeln
ist, wenn $n$ eine h\"ohere Potenz einer Primzahl ist,
immer gleich Null.}

Zur allgemeinen Bildung von $X_n$ wollen wir ein recurrentes
Verfahren anwenden; wir nehmen $X_n$ als schon gebildet an,
bezeichnen mit $p$ eine in $n$ nicht aufgehende Primzahl, mit $p'$
wie oben die $(\pi-1)^{\text{te}}$ Potenz von $p$, und bilden nun $X_{np'p}$,
worin nat\"urlich $p'$ auch gleich 1 sein kann.

Bezeichnen wir mit $r$ die \gesperrt{primitiven} $n^{\text{ten}}$ Einheitswurzeln,
mit $\alpha$ jede Einheitswurzel des Grades $pp'$, und mit $\alpha'$ jede Einheitswurzel des Grades $p'$, so erh\"alt man die s\"ammtlichen primitiven Einheitswurzeln des Grades $npp'$, wenn man von den
s\"ammtlichen $r\alpha$ die $r\alpha'$ wegnimmt (\S.\,132, II). Die $r\alpha$ sind aber
die Wurzeln der Gleichung:
\[
X_n(x^{pp'}) = 0,
\]
weil die Gr\"ossen
\[
(r\alpha)^{pp'} = r^{pp'},
\]
von der Reihenfolge abgesehen, mit den $r$ selbst \"ubereinstimmen
(\S.\,132, IV). Ebenso sind die $r\alpha'$ die Wurzeln der Gleichung
\[
X_n(x^{p'}) = 0;
\]
und daraus ergiebt sich
\weq{9}{X_{npp'}(x) = \frac{X_n(x^{pp'})}{X_n(x^{p'})}.}
Hiervon macht auch der Werth $n=1$ keine Ausnahme,
wenn wir unter $X_1$ die Function $x-1$ verstehen. Danach
k\"onnen wir $X_n$ in allen F\"allen verh\"altnissm\"assig einfach bilden.
Wir betrachten einige besondere F\"alle.

Nehmen wir an, es enthalte $n$ nur zwei von einander verschiedene Primzahlen $p$, $q$ und sei also
\[
n = pp'qq',
\]
dann ergiebt die Formel (8) und (9):
\[
X_n = \frac{(x^n-1)\left(x^{\frac{n}{pq}}-1\right)}{\left(x^{\frac{n}{p}}-1\right)\left(x^{\frac{n}{q}}-1\right)},
\]
%CONT>>
